\chapter{Missing data}
\label{ch:missing-data}

\begin{quote}
\emph{``The data you don't have is often more important than the data you do.''}
--- Anonymous statistician
\end{quote}

% ============================================================
\section{Why data goes missing}

Missing data is not a nuisance --- it is information. The \emph{reason} data is missing determines how you should handle it. Three mechanisms were formalized by Rubin (1976):

\begin{definition}[Missing data mechanisms]
\begin{itemize}
  \item \textbf{MCAR} (Missing Completely At Random): missingness is unrelated to any variable, observed or unobserved. Example: a lab sample is accidentally dropped.
  \item \textbf{MAR} (Missing At Random): missingness depends on observed variables but not on the missing value itself. Example: younger patients skip a depression questionnaire.
  \item \textbf{MNAR} (Missing Not At Random): missingness depends on the unobserved value. Example: high-income individuals refuse to report income.
\end{itemize}
\end{definition}

\begin{warningbox}
You cannot statistically prove MNAR --- it requires domain knowledge. But the distinction matters enormously: simple imputation works for MCAR, is reasonable for MAR, and can be badly biased for MNAR.
\end{warningbox}

% ============================================================
\section{Detecting missing values}

\begin{minted}{python}
import pandas as pd
import numpy as np

# Load Titanic data
url = ("https://raw.githubusercontent.com/datasciencedojo/"
       "datasets/master/titanic.csv")
df = pd.read_csv(url)

# Count missing values
print(df.isnull().sum())
print()
print(f"Total missing cells: {df.isnull().sum().sum()}")
print(f"Percentage missing: {df.isnull().mean().mean() * 100:.1f}%")
\end{minted}

\begin{minted}{python}
# Detailed missing report
def missing_report(df):
    """Return a DataFrame summarizing missing data."""
    missing = df.isnull().sum()
    pct = (missing / len(df)) * 100
    report = pd.DataFrame({"missing": missing, "pct": pct.round(1)})
    return report[report["missing"] > 0].sort_values("pct", ascending=False)

print(missing_report(df))
\end{minted}

\subsection{Hidden missing values}

\begin{minted}{python}
# Missing values are not always NaN
# Common disguises: "?", "N/A", "", -1, 999, -999, "unknown"
df_adult = pd.read_csv(
    "https://archive.ics.uci.edu/ml/machine-learning-databases/adult/adult.data",
    header=None,
    na_values=["?", " ?"],  # Adult dataset uses " ?" for missing
    names=["age", "workclass", "fnlwgt", "education", "education_num",
           "marital_status", "occupation", "relationship", "race",
           "sex", "capital_gain", "capital_loss", "hours_per_week",
           "native_country", "income"]
)
print(missing_report(df_adult))
\end{minted}

\begin{datatip}
Always scan for sentinel values: \texttt{-1}, \texttt{999}, \texttt{0} in numeric columns; empty strings, ``unknown'', ``N/A'' in text columns. Replace them with \texttt{np.nan} before any analysis.
\end{datatip}

% ============================================================
\section{Visualizing missing data}

\begin{minted}{python}
import missingno as msno
import matplotlib.pyplot as plt

# Matrix plot: white lines = missing
msno.matrix(df, figsize=(10, 5), sparkline=False)
plt.title("Missing data pattern — Titanic")
plt.tight_layout()
plt.savefig("missing_matrix.png", dpi=150)
plt.show()
\end{minted}

\begin{minted}{python}
# Bar chart: count of non-null values per column
msno.bar(df, figsize=(10, 4))
plt.tight_layout()
plt.show()

# Heatmap: correlation of missingness between columns
# Values near +1 mean columns tend to be missing together
msno.heatmap(df, figsize=(8, 6))
plt.tight_layout()
plt.show()
\end{minted}

\begin{minted}{python}
# Dendrogram: hierarchical clustering of missingness patterns
msno.dendrogram(df, figsize=(10, 5))
plt.tight_layout()
plt.show()
\end{minted}

\begin{preprocesstip}
The missingness correlation heatmap is extremely useful: if two columns have correlated missingness, they may share a common cause (e.g., both come from the same survey section that some respondents skipped).
\end{preprocesstip}

% ============================================================
\section{Deletion strategies}

\subsection{Listwise deletion (dropping rows)}

\begin{minted}{python}
# Drop rows with ANY missing value
df_complete = df.dropna()
print(f"Before: {len(df)}, After: {len(df_complete)}")
print(f"Lost: {len(df) - len(df_complete)} rows "
      f"({(len(df) - len(df_complete)) / len(df) * 100:.1f}%)")
\end{minted}

\subsection{Column deletion}

\begin{minted}{python}
# Drop columns with more than 50% missing
threshold = 0.5
cols_to_drop = df.columns[df.isnull().mean() > threshold].tolist()
print(f"Dropping columns: {cols_to_drop}")
df_reduced = df.drop(columns=cols_to_drop)
\end{minted}

\begin{warningbox}
Listwise deletion is only safe under MCAR. Under MAR or MNAR, it introduces bias. For the Titanic dataset, older passengers with higher-class tickets are more likely to have recorded ages --- dropping missing ages skews the dataset toward lower-class passengers.
\end{warningbox}

% ============================================================
\section{Simple imputation}

\subsection{Mean, median, and mode imputation}

\begin{minted}{python}
from sklearn.impute import SimpleImputer

# Numerical: median imputation (robust to outliers)
num_imputer = SimpleImputer(strategy="median")
df["Age"] = num_imputer.fit_transform(df[["Age"]])

# Categorical: mode (most frequent) imputation
cat_imputer = SimpleImputer(strategy="most_frequent")
df["Embarked"] = cat_imputer.fit_transform(df[["Embarked"]]).ravel()

print(df.isnull().sum())
\end{minted}

\subsection{Constant imputation}

\begin{minted}{python}
# Fill with a specific value
df["Cabin"] = df["Cabin"].fillna("Unknown")

# Fill numerical with 0 (use with caution)
df["some_col"] = df["some_col"].fillna(0)
\end{minted}

% ============================================================
\section{Advanced imputation}

\subsection{KNN imputation}

\begin{minted}{python}
from sklearn.impute import KNNImputer

# KNN uses similar rows to estimate missing values
# Only works on numerical data
num_cols = ["Age", "Fare", "SibSp", "Parch"]
knn_imputer = KNNImputer(n_neighbors=5, weights="distance")
df[num_cols] = knn_imputer.fit_transform(df[num_cols])

print(f"Missing after KNN: {df[num_cols].isnull().sum().sum()}")
\end{minted}

\subsection{Iterative imputation (MICE)}

\begin{minted}{python}
from sklearn.experimental import enable_iterative_imputer
from sklearn.impute import IterativeImputer

# Iterative imputer models each feature as a function of others
iter_imputer = IterativeImputer(max_iter=10, random_state=42)
df[num_cols] = iter_imputer.fit_transform(df[num_cols])

print("Iterative imputation complete.")
print(df[num_cols].describe())
\end{minted}

\begin{preprocesstip}
KNN imputation respects local structure (similar passengers get similar ages). Iterative imputation captures linear relationships across features. Both outperform mean/median imputation when data is MAR.
\end{preprocesstip}

% ============================================================
\section{Imputation with an indicator}

\begin{minted}{python}
# Create a binary indicator BEFORE imputing
df["Age_was_missing"] = df["Age"].isnull().astype(int)

# Then impute
df["Age"] = df["Age"].fillna(df["Age"].median())

# The indicator preserves the information that the value was missing
print(df[["Age", "Age_was_missing"]].head(10))
\end{minted}

\begin{datatip}
Adding a missing indicator column is one of the most underused techniques. It lets downstream models learn whether the missingness itself is predictive. In the Titanic dataset, having a missing cabin number correlates with lower survival.
\end{datatip}

% ============================================================
\section{Choosing an imputation strategy}

\begin{longtable}{p{3.5cm}p{3cm}p{3cm}p{3.5cm}}
\toprule
\textbf{Method} & \textbf{Best for} & \textbf{Weakness} & \textbf{Mechanism} \\
\midrule
Drop rows & $<$5\% missing & Loses data & MCAR only \\
Mean/Median & Quick baseline & Reduces variance & MCAR \\
Mode & Categorical & Ignores structure & MCAR \\
KNN & Clustered data & Slow on large data & MAR \\
Iterative (MICE) & Linear relationships & Assumes linearity & MAR \\
Domain constant & Known meaning & May introduce bias & Any \\
\bottomrule
\end{longtable}

% ============================================================
\section{Exercises}

\begin{exercise}
\begin{enumerate}
  \item Load the Melbourne Housing dataset. Create a missing data report and a \texttt{missingno} matrix plot. Which columns have correlated missingness?
  \item For the Titanic dataset, compare the distribution of \texttt{Age} before and after mean imputation vs.\ KNN imputation. Plot both distributions on the same histogram.
  \item Load the Adult Census dataset (with \texttt{na\_values=[" ?"]}) and impute the three categorical columns with missing values using mode imputation. Verify no missing values remain.
  \item Implement a function \texttt{smart\_impute(df)} that: (a) drops columns with $>$70\% missing, (b) imputes numerical columns with median, (c) imputes categorical columns with mode, (d) adds missing indicators for all originally-missing columns.
  \item Investigate whether \texttt{Age} in the Titanic dataset is MCAR by comparing the survival rate of passengers with and without a recorded age. What do you conclude?
\end{enumerate}
\end{exercise}

% ============================================================
\section{Chapter summary}

\begin{itemize}
  \item Missing data has three mechanisms: MCAR, MAR, and MNAR. The mechanism determines which imputation strategy is valid.
  \item Always scan for hidden missing values (sentinel values like \texttt{?}, \texttt{-1}, \texttt{999}).
  \item The \texttt{missingno} library provides matrix, bar, heatmap, and dendrogram visualizations for missing patterns.
  \item Simple imputation (mean, median, mode) is fast but reduces variance and only works well under MCAR.
  \item KNN and iterative (MICE) imputation capture relationships between features and work under MAR.
  \item Adding a missing indicator column preserves the information content of missingness.
\end{itemize}
